% STATUS: LITERATURE AUDIT COMPLETE WITH DISCLOSED ACCESS GAPS
% Phase 1B, 24 September 2026. Includable fragment; no project proofs.
\section{Phase 1B: adversarial novelty audit}
\label{sec:adversarial-audit}

\subsection{Verdict and limits of closure}
The original eight targets do not presently support eight novelty claims.
Four basic claims are unsafe as principal contributions: Fourier diagonalization,
the fixed-symbol large-network limit, Fourier inversion for static response,
and generic perturbation bounds. Mode selection and the first Lyapunov
coefficient require substantial sharpening. The matched comparison is a
potentially useful specialization. Priority for the demand-modulated coupling
law remains unresolved. No target receives GREEN in this audit.

This is a source-based assessment, not verification of the blueprint's mathematics.
The canonical demand-modulated model (DM), matched bilinear comparator (MB),
positive parameter domain, indexing, and terminology remain unchanged.
All P01--P18 files remain TODO. Neither an unsuccessful search nor a different
displayed equation establishes novelty. No proof or numerical experiment was
performed. Recommendations below are proposals for subsequent author decisions.

The six supplied PDFs and the original Phase 1 audit remain the baseline.
This revision adds English and Chinese searches, backward citation tracing,
and stronger generic counterexamples to the novelty narrative.
The search cutoff is 24 September 2026. Access restrictions prevent full
closure of several five-element sources. Their missing equations are recorded
as missing, not reconstructed from abstracts. A separate archival research
track is neither created nor recommended.

\subsection{Evidence convention and source locations}
S01--S14 retain their Phase 1 identifiers. Page numbers for S01--S04 and
S08 are those of the inspected author manuscripts, not silently converted
journal pages. New records specify their version. ``Proves'' below means
that an argument is present in the source, not that this audit independently
certifies every line. ``Stronger'' refers only to the stated overlapping
component and its assumptions; it does not mean that a different vector field
contains DM. Search-index excerpts of primary PDFs are explicitly distinguished
from full-text inspection. The files \texttt{additional\_sources\_needed.md}
and \texttt{audits/phase1b\_search\_log.md} document the unresolved frontier.

\subsubsection{Reassessment of the original six papers}
\begin{description}[style=nextline,leftmargin=0pt,labelwidth=0pt,labelsep=0pt]
\item[S01: Jaramillo, Mrad and Stepien (2023).]
\cite{JaramilloMradStepien2023}; DOI \href{https://doi.org/10.1063/5.0138150}{10.1063/5.0138150}.
In arXiv:2210.04342v2, \S III, p.~3, (5)--(6) define the
three-species May--Leonard system with linear mutation. \S IV.B, p.~5,
Theorem~1 and (7)--(8) treat its circulant spectrum; \S IV.C, pp.~6--7,
Theorem~2 and (15)--(25) give a Hopf/Lyapunov calculation, including Fourier
vectors in (22)--(23). The competition coefficients satisfy $\alpha,\beta>0$;
the stability discussion separates $\mu\geq1/6$ from $0\leq\mu<1/6$.
The Hopf calculation uses the critical line (11) and a nonzero frequency;
one must exclude its zero-frequency coincidence. The proofs concern that
three-species polynomial system. Relation: orthogonal model, overlapping
Hopf calculation method; not an arbitrary-$(n,q)$ DM theorem.

\item[S02: Pais, Caicedo-N\'u\~nez and Leonard (2012).]
\cite{PaisCaicedoLeonard2012}; DOI \href{https://doi.org/10.1137/120878537}{10.1137/120878537}.
The supplied 30-page manuscript uses a replicator--mutator model on a simplex,
\S2, (2.2), pp.~3--5. The circulant family in \S4, (4.1), p.~10,
has $N\geq3$, the specified matrix entries and mutation structure; its
parameters satisfy $\alpha,\beta\in[0,1)$ and $\alpha+\beta>0$.
\S4.1, pp.~11--13, (4.2)--(4.4), Lemmas~4.2/4.4 and Theorem~4.5
give mode/Hopf results. Lemma~4.6, p.~13, and Appendix~D,
pp.~26--28, (D.1)--(D.4), reduce the Lyapunov coefficient by Fourier
structure, including cubic terms. Appendix~B, pp.~24--25, records
the general formula. Relation: different state space and kinetics;
strong direct precedent for the proposed method, including the cubic part.

\item[S03: van Kekem and Sterk (2018): the closest theorem-program precedent.]
\cite{vanKekemSterk2018}; DOI
\href{https://doi.org/10.1016/j.physd.2017.11.008}{10.1016/j.physd.2017.11.008}.
In the 55-page accepted manuscript, (1), p.~1, is Lorenz--96, with
scalar forcing $F$ and a homogeneous equilibrium. \S2.1, pp.~8--9,
(3)--(6) gives a circulant spectrum with real dispersion involving
$\cos\theta-\cos2\theta$. Proposition~5, pp.~11--12, selects the first
unstable mode for $F>0$, $n\geq4$, maximizing over $0<l<n/3$;
its bounds have an $n=7$ exception and $n=12$ has a two-pair tie.
Theorem~2, p.~10, treats simultaneous crossings. Theorem~4, p.~11,
gives an explicit first Lyapunov coefficient and criticality information
for $n\geq4$, $0<l<n/2$, $l\neq n/3$, when that pair is the only
one crossing at $F_H(l,n)$.
Proposition~6, \S2.3, p.~13, gives the selected-wave limit
$n/l_1(n)\to 2\pi/\arccos(1/4)$, with its argument on p.~13.
The Lyapunov/selection proofs appear in \S\S4.3--4.4, pp.~38--48.
These are proved results for Lorenz--96, not DM.
Relation: orthogonal equation, but substantially the same proposed
sequence of spectral selection, large-$n$ asymptotics and Hopf calculation.
Changing coefficients in a comparable dispersion relation is insufficient
by itself to establish a substantive new theorem program.

\item[S04: Wei et al. (2021).]
\cite{WeiEtAl2021}; DOI \href{https://doi.org/10.1016/j.jfranklin.2020.11.011}{10.1016/j.jfranklin.2020.11.011}.
The supplied 18-page manuscript defines the five-dimensional dynamo at
(2.1), p.~3; Theorems~2.1--2.2, pp.~4--5, discuss Hopf points using
reported numerical quantities. Theorems~3.1--3.2, \S3, pp.~5--7,
concern the specified controlled/observer error system. Five coordinates
do not make this a scalar circulant ring. Relation: orthogonal to the
DM coupling law and passive parameter-robustness bounds. Neither its
control conclusions nor its reported decimals establish our claims.

\item[S05: Lin and Han (2020): exact comparator overlap at five nodes.]
\cite{LinHan2020}; DOI \href{https://doi.org/10.23919/CCC50068.2020.9188438}{10.23919/CCC50068.2020.9188438}.
\S2.1, p.~778, (2) is the six-term homogeneous five-element polynomial
model with nonnegative coefficients. Our $\MB_{5,2}$ is its exact
subfamily under the coefficient assignment
$(a_L,b_L,c_L,d_L,g_L,h_L)=(a,k,c,h,0,g)$.
\S2.2, p.~778, discusses positivity; \S2.2.1, Case~4, p.~779,
(5)--(7), addresses homogeneous equilibria. Numerical examples are in
\S\S2.2.2--2.2.3, pp.~779--781. Relation: the source model is
stronger in parameter count at $n=5$; MB is identical on that subfamily.
This is a source-equation identification, not a new analysis of MB.
The comparison with DM still requires a distinct result. References on
p.~781 lead to earlier Chinese and English five-element models.

\item[S06: Nian and Li (2025).]
\cite{NianLi2025}; DOI \href{https://doi.org/10.1186/s13662-025-03909-9}{10.1186/s13662-025-03909-9}.
In the publisher PDF, \S2.2, pp.~5--6, (1) defines a heterogeneous
five-variable linear/quartic model; (3), p.~7, is its Jacobian.
\S\S3--4, pp.~8--13, report fitted parameters and dynamical computations.
Relation: orthogonal coupling law, overlapping five-element motivation.
Reported simulations do not prove arbitrary-ring statements or validate
the biological interpretation of DM.
\end{description}

\subsubsection{Generic theory that defeats broad novelty claims}
\begin{description}[style=nextline,leftmargin=0pt,labelwidth=0pt,labelsep=0pt]
\item[S07: Gray (2006).]
\cite{Gray2006}; \href{https://ee.stanford.edu/~gray/toeplitz.pdf}{author PDF},
\S3.1, printed pp.~32--34, (3.7)--(3.13); \S3.2,
Theorem~3.1, printed pp.~34--35. Finite circulant matrices have a unitary
Fourier decomposition; inverse formulas require nonzero eigenvalues.
These results are established in the source. Relation: strictly more
general linear algebra than our homogeneous Jacobian, once that Jacobian
has been correctly derived. Spectrum, normality and inverse Fourier sums
are supporting applications, not new general theorems.

\item[S08: Kerin and Engler, now identified in its published form (2022).]
\cite{KerinEngler2022}; DOI \href{https://doi.org/10.3934/dcdsb.2021064}{10.3934/dcdsb.2021064}.
The inspected version remains arXiv:2005.07767v2: \S3.2,
pp.~11--12, (27)--(29), Lemma~1; Appendix~A, p.~27,
Proposition~4, (48)--(49). Assumptions include a circulant linear part,
an equivariant quadratic nonlinearity and a simple nonresonant Hopf
setting with the displayed inverses defined. The scalar Fourier
reduction is proved for this class. Relation: broader quadratic class,
but not the full rational DM nonlinearity. \S5, pp.~21--24,
(40), Propositions~1--3, allows inhomogeneous coefficients under
energy-preserving quadratic and positive damping hypotheses; those
hypotheses must not be imported into DM. The new bibliography entry
adds the journal record while preserving the original preprint entry.

\item[S09 and S22: equivariant Hopf and phase patterns.]
Stewart \cite{Stewart2024}, arXiv:2403.15814v1,
\S3, Theorem~3.7, and \S\S5.2--5.4, (16)--(17),
Theorems~5.5/5.7 and Corollary~5.6, supplies a ring-specific synthesis.
The original research source Filipski--Golubitsky
\cite{FilipskiGolubitsky2010}, DOI
\href{https://doi.org/10.1137/090756582}{10.1137/090756582},
\S3, p.~285, Theorem~3.1, and \S5, p.~288, Theorem~5.1,
proves abelian equivariant Hopf symmetry restrictions/realization.
The hypotheses concern smooth generic equivariant bifurcation, a critical
representation and isotropy/kernel conditions; the realization theorem
is not existence in every particular vector field. Relation: general
symmetry theory subsumes our conditional rotating-wave conclusion.
DM must still supply the simple imaginary pair, crossing, nondegeneracy
and all-mode stability conditions. Do not copy the phase expression in
Stewart v1 (15) without reconciling it with the surrounding convention.

\item[S10--S11: matrix perturbation and threshold sensitivity.]
Embree--Trefethen \cite{EmbreeTrefethen2001}, DOI
\href{https://doi.org/10.1137/S1064827500373012}{10.1137/S1064827500373012},
\S1, pp.~583--584, Theorem~4$\varepsilon$, p.~585, and the normal-matrix
discussion on p.~586, relates Euclidean perturbations, resolvents and
pseudospectra. Its normal-matrix case is broader than our stable
homogeneous linearization. Greenbaum--Li--Overton
\cite{GreenbaumLiOverton2020}, DOI
\href{https://doi.org/10.1137/19M124784X}{10.1137/19M124784X},
author PDF \S2, p.~4, Assumption~1, Theorem~1 and (1), proves
the derivative of a simple eigenvalue for an analytic matrix family
with normalized left/right eigenvectors. Relation: the general derivative
is stronger than its DM application. Differentiating a simple transverse
threshold is routine after the moving equilibrium has been included.

\item[S19: Benzi and Razouk (2007), a stronger localization baseline.]
\cite{BenziRazouk2007};
\href{https://www.cs.emory.edu/~benzi/Web_papers/benzi_razouk.pdf}{published ETNA PDF}.
\S3.5, p.~24, Theorem~3.4 and (3.7), proves decay in directed graph
distance for diagonalizable $A=XDX^{-1}$ whose spectrum lies in a
continuum $F$, with $f$ analytic on the larger domain $I(C_{R_0})$.
The bound contains $\kappa(X)$. \S3.7, p.~27, Theorem~3.7 and
(3.10), supplies explicit constants under additional geometric hypotheses.
Relation: stronger generic matrix-function theory; using $f(z)=1/z$
requires the analytic region to avoid zero. Uniform-in-$n$ claims need
uniform hypotheses, not just invertibility at each size. A cyclic graph
distance must respect wraparound; ordinary band distance need not.
This closes the generic decay-method gap even though the full S12
Demko--Moss--Smith text \cite{DemkoMossSmith1984} remains only partly accessible.

\item[S20: Sun, Bollt and Nishikawa (2009).]
\cite{SunBolltNishikawa2009}; DOI
\href{https://doi.org/10.1209/0295-5075/85/60011}{10.1209/0295-5075/85/60011}.
In arXiv:0811.0649v1, p.~2, (1) is a uniformly coupled
near-identical network; pp.~4--5, (6)--(9) diagonalize its
inhomogeneous variational equation when the modified Laplacian is
diagonalizable. Pp.~6--7, (12)--(15), derive a forced-response bound
under an exponential transition bound and bounded mismatch forcing.
These are linearized conclusions; replacement of the mean trajectory by
an isolated unit is explicitly conjectural on p.~6. Relation: overlapping
method, different coupling class. It does not prove nonlinear DM
persistence or a global robustness radius.

\item[S21: Qiu et al. (1995).]
\cite{QiuEtAl1995}; DOI
\href{https://doi.org/10.1016/0005-1098(95)00024-Q}{10.1016/0005-1098(95)00024-Q}.
\S1, pp.~879--880, (1)--(3), treats stable $A+BDC$ with fixed
real $A,B,C$ and a full real or complex perturbation block $D$ measured
in the spectral norm. The unnumbered Theorem, p.~886, (20), is the
real structured-singular-value formula, with proof in \S2.
Relation: stronger exact linear robustness theory, but its full-block
uncertainty is not arbitrary repeated/diagonal parameter uncertainty.
It does not include a nonlinear equilibrium's displacement automatically.
Backward reference: Hinrichsen--Pritchard (1986), DOI
\href{https://doi.org/10.1016/0167-6911(86)90094-0}{10.1016/0167-6911(86)90094-0};
its full text was not inspected here.

\item[S23: Asllani et al. (2014).]
\cite{AsllaniEtAl2014}; DOI
\href{https://doi.org/10.1038/ncomms5517}{10.1038/ncomms5517}.
In the 28-page arXiv:1402.0760 manuscript, \S II,
pp.~8--10, (6)--(11), develops directed-network dispersion and an
instability condition. Assumptions include a stable two-species reactive
equilibrium and a complete Laplacian eigenbasis; p.~8, footnote~2,
requires balance for the Fickian transport convention to preserve
homogeneity. The linear instability analysis is analytic; subsequent
travelling patterns are illustrated numerically. Relation: broader graph
setting but different multi-species diffusion kinetics, not exact
subsumption of DM selection. Directed dispersion and pattern formation
are already established subjects.

\item[S24: Montanari, Zanin and Motter (2026): recent priority warning.]
\cite{MontanariZaninMotter2026};
\href{https://arxiv.org/html/2609.25226v1}{arXiv:2609.25226v1},
submitted 21 September 2026. Supplement \S S3, Assumptions~1--2,
requires diagonalizable $L$, $N\geq3$, nonnegative real parts of its
eigenvalues, and nonnegative damping. \S S3.2, Theorem~4,
for (3)--(4), assumes non-Hermitian circulant $L$ and differentiability
at the homogeneous optimum; it proves this optimum is not locally
optimal among heterogeneous damping choices. Equations (S49)--(S54)
use second derivatives and Fourier mode sums. \S S4 separates proved
results from a mode-mixing conjecture. Relation: different second-order
network model, substantial overlap with broad heterogeneity-improves-stability
or mode-mixing narratives. This preprint has not been independently
validated and is not evidence that every perturbation stabilizes DM.
\end{description}

\subsubsection{Five-element and biochemical priority search}
\begin{description}[style=nextline,leftmargin=0pt,labelwidth=0pt,labelsep=0pt]
\item[S13--S16: the Lin--Shi--Han enzymatic sequence is the main access gap.]
S13 \cite{LinShiHan2021}, DOI
\href{https://doi.org/10.23919/CCC52363.2021.9550151}{10.23919/CCC52363.2021.9550151},
has pages 6410--6415, now confirmed by S14 reference~37.
S15 \cite{LinShiHan2022}, DOI
\href{https://doi.org/10.23919/CCC55666.2022.9901543}{10.23919/CCC55666.2022.9901543},
pp.~5735--5740, and S16 \cite{LinHan2023}, DOI
\href{https://doi.org/10.23919/CCC58697.2023.10239797}{10.23919/CCC58697.2023.10239797},
pp.~6823--6828, are identified in S14 references~38--39.
S14 \cite{LinHan2024}, DOI
\href{https://doi.org/10.1007/s11071-024-09374-5}{10.1007/s11071-024-09374-5},
publisher abstract and references, studies enzymatic regulatory network
topologies and the role of five nodes. Complete equations/supplement
were not retrieved. Exact assumptions, equation identity, and theorem
overlap therefore remain GRAY for this sequence. An indexed author
upload under the 2022 title contains a later uncorrected proof: versions
must not be conflated. No theorem-level adverse or favorable priority
conclusion is based on that upload.

\item[S17: Zhang, Xia and Wang (2020).]
\cite{ZhangXiaWang2020}; DOI
\href{https://doi.org/10.1088/1674-1056/abc2bf}{10.1088/1674-1056/abc2bf}.
Publisher-PDF indexed text at \S5.1, p.~120504-5, (1), specifies
five nodes with production formed from Hill activation and Hill
repression factors, and linear degradation. The Hill exponents are
not network size. This is a definite prior rational/saturating Wu-Xing
coupling construction; it differs from the displayed DM law.
The publisher abstract reports a landscape/flux investigation.
Access is partial: the equation passage and metadata were inspected,
not the full numerical argument. Relation: orthogonal exact law,
direct overlap with any broad claim to introduce saturation into
five-element networks. Reference~11, p.~120504-11, leads to
Yang et al. (1995), requiring retrieval.

\item[S18: Ma et al. (2009), a backward biochemical source.]
\cite{MaEtAl2009}; DOI
\href{https://doi.org/10.1016/j.cell.2009.06.013}{10.1016/j.cell.2009.06.013}.
Experimental Procedures, ``Equations of the Circuit'', p.~772, (7),
uses enzyme activities normalized to fixed total concentration one,
positive kinetic constants, and source activity multiplied by saturating
recipient activation/deactivation factors. Its three-node topology
search and associated adaptation analysis concern that biochemical
class. Relation: recipient-state modulation is already a modeling
ingredient; the factors are not the DM fractions. This prevents
equating a new name with a new coupling principle. A coordinate or
parameter equivalence to DM would need actual verification.

\item[S25: Ganikhodjaev, Pah and Rozikov (2019).]
\cite{GanikhodjaevPahRozikov2019};
\href{https://arxiv.org/pdf/1908.07740}{arXiv:1908.07740}.
The model is a discrete quadratic stochastic operator on a simplex,
not an ODE. \S4.3, p.~11, Theorem~1, and \S4.5,
pp.~13--14, Proposition~7/Theorem~4, contain dynamical results for
the specified parameter/operator class. Relation: orthogonal to DM
and continuous-time MB; relevant to prior mathematical use of the
five-element graph, not evidence for an identical differential law.

\item[S26--S29 and older Chinese sources: explicitly unresolved.]
Ding--Wang (2008) \cite{DingWang2008}, DOI
\href{https://doi.org/10.1007/s11726-008-0269-8}{10.1007/s11726-008-0269-8},
pp.~269--270: publisher abstract/reference list inspected, equations
unavailable. Its references lead to Li (1987)
\cite{Li1987}, DOI
\href{https://doi.org/10.1088/0031-8949/36/6/015}{10.1088/0031-8949/36/6/015},
pp.~966--969, and Zhang--Han (1998), pp.~130--135.
Lin--Han (2020), p.~781, also points to Feng--Niu--Wang (2003),
pp.~71--74, Wang's 2008 thesis, Bellavite et al. (1998)
\cite{BellaviteEtAl1998}, and Kim et al. (2014)
\cite{KimEtAl2014}. Abstracts/metadata are insufficient to exclude
transformed rational or bilinear models. No missing theorem number
or equation locator is invented. The retrieval list states exactly
what must be inspected before T1 can be cleared.
\end{description}

\subsection{What a conditional Hopf theorem would and would not add}
Known theory supplies the local bifurcation mechanism and permissible
spatiotemporal symmetry. Fourier mode arithmetic and use of the zero
and second harmonics in a Lyapunov calculation also have explicit
precedents (S02, S08). These are not separate discoveries.
The still-unproved DM content is an admissible positive equilibrium
branch, the first crossing on a specified path, exclusion of simultaneous
critical modes and zero frequency, transversality including equilibrium
motion, and a rigorously certified nonzero Lyapunov coefficient.
An open parameter region with a qualitative MB contrast would make these
checks serve one mathematical question. A list of conditional formulas
with no certified realizability would not.

Neither a maximum real part at one parameter value nor a tied pair of
conjugate modes proves a first Hopf bifurcation. Distinct critical pairs,
the real Nyquist mode when $n$ is even, and modular harmonic coincidences
must be separated. Equivariant theory does not supply orbital stability
without the relevant remaining spectrum and criticality conditions.
Exact nonlinear phase symmetry must be distinguished from a leading
sinusoidal approximation; its sign depends on the frozen shift convention.

\subsection{Assessment of the four proposed packages}
\begin{description}[style=nextline,leftmargin=0pt,labelwidth=0pt,labelsep=0pt]
\item[A: spectrum + step-two selection + large-network limit.]
Likely significance in its basic form: modest. Diagonalization and
continuous/discrete maximization are largely generic; S03 is a particularly
close precedent. The coefficient-to-parameter map and realizable mode
regions would be model-specific. Referee objection: a quadratic optimization
exercise appended to a new equation. Needed addition: a complete admissible
parameter classification tied to a qualitative dynamical/comparator contrast,
with exceptional sizes and ties treated. A limit alone does not cure this.

\item[B: selection + rigorous Hopf + rotating-wave branch.]
Potentially coherent, but generic Hopf and equivariance constitute much
of the machinery. The useful specific content would be a certified first
instability and criticality phase diagram on positive branches.
Referee objection: S02/S03 already perform an analogous program.
Needed addition: analytic realizability on nonempty open parameter regions,
a meaningful MB contrast, and rigorous criticality. Merely finding an
imaginary pair and invoking symmetry is insufficient.

\item[C: selection + explicit Lyapunov coefficient + structured heterogeneity.]
Potentially stronger if the last component explains a new mechanism.
The tensor reduction, inverse Fourier sum, and first-order sensitivity
are mostly known. DM derivatives and sharp parameter-dependent response
could be specific. Referee objection: three routine calculations with
little interaction. Needed addition: a sharp structured result tied to
the selected mode or its nonlinear branch, with a controlled nonlinear
remainder and a quantified gain over the generic baselines S19/S21.
An assertion that disorder can help is insufficient in light of S24.

\item[D: all of B and C.]
Largest scope, not automatically greatest significance. The model-specific
part would have to link wave selection, branch criticality and a meaningful
heterogeneity mechanism. Referee objection: overextended catalogue with
the difficult assertions left conditional. Needed addition: one unifying
theorem and a clear reason the heterogeneity result belongs to the main
argument. More standard lemmas do not replace such a result.
\end{description}

\paragraph{Recommended smallest package.}
Pursue a strengthened B with P13 criticality and P08 comparison as
essential parts: an exact, realizable $q=2$ selection/first-Hopf
classification that explains a qualitative difference between DM and MB.
Keep P04 as a supporting general lemma, P06 as a useful quantitative
corollary only if it clarifies the mechanism, and P14 as the symmetry
conclusion. Do not make the full P15--P18 extension a prerequisite for
this core. If no such model-specific contrast or sign classification
survives subsequent work, reconsider the contribution before expanding
the theorem count. This recommendation is conditional, not a prediction
of journal acceptance or a claim that the proposed strengthened results
are true.

\subsection{Recommended scope of the generalization}
Retain arbitrary $(n,q)$ for the model definition and general spectral
background. Make $q=2$ the substantive setting. At the level of a
candidate symbol using the first and second spatial harmonics, the
standard identity for $\cos2\theta$ explains the quadratic structure;
this observation is not a derivation of the DM coefficients.
For other $q$, the algebra and competing modes need separate treatment.
No complete arbitrary-$q$ selection theorem is recommended now.

Varying $n$ permits more mode classes, parity effects and possible
mode competition unavailable at five nodes. These are motivations for
the generalization, not verified DM phenomena. Show that the admissible
parameters actually realize relevant changes; relabeling a sampled
maximizer is not enough. Large-$n$ work is useful when it predicts wavelength
and finite-size threshold error under a fixed, explicit family. It is
weak if it only replaces a grid by a continuum or changes parameters
with $n$ without uniform control.

Present $n=5,q=2$ as the motivating example and a checkable specialization,
with prior Wu-Xing models credited in the main audit. Golden-ratio
trigonometric constants are not intrinsic novelty. The condition
$\gcd(n,q)=1$ concerns the step-$q$ subcycle; it is not necessary for
the full network's connectivity when the step-one edge is present.
No restriction of the canonical parameter domain or alteration of the
coupling law is justified solely by this novelty audit.

\subsection{Closure gates before claiming a substantive contribution}
The eight target decisions are in the accompanying revised matrix.
Module dispositions and precise proposed strengthenings are in
\texttt{novelty\_decisions.md}. Priority for T1 remains blocked by the
enzymatic sequence and older source equations. This does not prevent
later independent mathematics if explicitly authorized; it prevents a
responsible priority claim. Branch/path, comparison calibration,
heterogeneity assignment and uncertainty norm also remain open choices.
No result has been promoted to VERIFIED or AUDITED by this literature work.
